\item[\textbf{5.}] \textbf{[30 min/20 puntos]} \textcolor{dkgreen}{Peterson es una implementación de un candado para dos hilos. En la implementación siguiente (\textit{DoublePeterson}) se utilizan tres candados Peterson para realizar una implementación de un candado para 4 hilos; se asume que a cada hilo le corresponde un ID único entre $1$ y $4$.}

    \textcolor{dkgreen}{Recuerda que en el candado de Peterson se cumplen las propiedades de Exclusión mutua, Deadlock-free y Starvation-free. Argumenta si \textit{DoublePeterson} cumple con:}

    \begin{enumerate}
        \item \textcolor{dkgreen}{Exclusión mutua.}
        \item \textcolor{dkgreen}{Deadlock-free.}
        \item \textcolor{dkgreen}{Starvation-free.}
        \item \textcolor{dkgreen}{Justicia.}
    \end{enumerate}

    \textcolor{red}{\textbf{Hint:} Si cumple con la propiedad, muestra la invariante que se mantiene (en el caso de la Exclusión mutua y Justicia), o muestra la cosa buena que pasará (en el caso de Deadlock-free y Starvation-free). No hace falta que demuestres Peterson.}

    \textcolor{dkgreen}{Si no cumple con alguna propiedad, muestra una ejecución que lo justifique.}

    \begin{lstlisting}[language=Java]
class DoublePeterson implements Lock { // Candado para n = 4 hilos
    Shared Variables:
        1: Lock lock1, lock2, lock3;
        2: DoublePeterson() {
        3:     lock1 = petersonLock();
        4:     lock2 = petersonLock();
        5:     lock3 = petersonLock();
        6: }

        7: void lock() {
        8:     int me = ThreadID.get(); // Cada hilo tiene un ID único de 1 a 4
        9:     if (me == 1 or me == 2) {
        10:        lock1.lock(); // Los hilos 1 y 2 toman el candado 1
        11:    }
        12:    if (me == 3 or me == 4) {
        13:        lock2.lock(); // Los hilos 3 y 4 toman el candado 2
        14:    }
        15:    lock3.lock(); // Todos pueden tomar el candado 3
        16: }

        17: void unlock() {
        18:     int me = ThreadID.get();
        19:     lock3.unlock();
        20:     if (me == 1 or me == 2) {
        21:        lock1.unlock();
        22:    }
        23:    if (me == 3 or me == 4) {
        24:        lock2.unlock();
        25:    }
        26: }
}
    \end{lstlisting}

    \begin{center}
        \textbf{Figura 2:} Implementación de un candado ``DoublePeterson''.
    \end{center}
\begin{enumerate}
    \item Exclusión Mutua \\
    La invariante sería, que para que cualquier hilo logre entrar a la sección crítica, debe haber ejecutado exitosamente \texttt{lock3.lock()}. Como \texttt{lock3} es un candado de Peterson, su invariante garantiza que a lo más, un solo hilo puede poseerlo a la vez. Al ser \texttt{lock3} el último paso del método \texttt{lock()} del \texttt{DoublePeterson}, es imposible que dos hilos esten en la CS simultaneamente. Por lo tanto, si cumple con Exclusión Mutua.
    \item Deadlock-free \\
    Si varios hilos intentan entrar, el sistema siempre avanzará. Sabemos que los candados de Peterson son deadlock-free. Si los hilos 1 y 2 compiten en \texttt{lock1}, al menos uno ganará. Si los hilos 3 y 4 compiten en \texttt{lock2}, al menos uno ganará. A lo más, habrá dos ganadores compitiendo por \texttt{lock3}. Como \texttt{lock3} también es deadlock-free, obligatoriamente se resolverá el empate y dejará pasar a uno. El sistema nunca se congela. Por lo tanto \texttt{DoublePeterson} es deadlock-free.
    \item Starvation-free \\
    Si un hilo particular (digamos el Hilo 1) quiere entrar a la CS, eventualmente lo logrará. Primero competirá por \texttt{lock1}, y como \texttt{Peterson} es starvation-free, el hilo 2 no puede permanecer ahí infinitamente, así que el hilo 1 eventualmente ganará \texttt{lock1}. Una vez con \texttt{lock1}, el hilo 1 competirá por \texttt{lock3}, en la que nuevamente este tiene la propiedad de starvation-free, por lo que su entrada eventual a la CS está garnatizada. Por lo tanto, \texttt{DoublePeterson} cumple starvation-free.
    \item FIFO (Justicia) \\
    Supongamos que se inició el doorway en todos lo hilos del sistema.
    \begin{enumerate}
        \item El hilo 1 quiere entrar. Ejecuta \texttt{lock1.lock()}. Realiza su doorway completo (levanta su bandera y se hace víctima en el candado 1). En ese momento, el sistema lo pausa, justo antes de evaluar su \texttt{while}, es decir \texttt{D_1} ya terminó.
        \item El hilo 3 quiere entrar. Ejecuta \texttt{lock2.lock()}. El hilo 3 pasa su doorway, y no se detiene en el ciclo \texttt{while}, ganando \texttt{lock2}.
        \item El hilo 3 avanza y ejecuta \texttt{lock3.lock()}. Como nadie más está compitiendo al final todavía, pasa el doorway y no se detiene en el \texttt{while} y entra a la sección crítica.
    \end{enumerate}
    Entonces, el hilo 1 terminó su doorway global antes de que el hilo 3 empezara a perdir el candado ($D_1 \rightarrow D_3$), pero el hilo 3 logró entrar a la CS primero ($CS_3 \rightarrow CS_1$). Por lo tanto no cumple FIFO.
\end{enumerate}